% unit: noether.p09.title.001
\begin{center}
{\Large\bfseries 9. Les domaines les plus généraux d'entiers transcendants}\par
\vspace{0.75em}
Par Emmy Noether, à Erlangen.\par
\vspace{1.2em}
\emph{Math. Ann.} 77 (1916), p. 103--128
\end{center}

\vspace{1.5em}

% unit: noether.p09.par.001
Au moyen du théorème du bon ordre, E. Zermelo a construit un domaine d'\og{} entiers transcendants \fg{}, c'est-à-dire un domaine numérique $\mathfrak G$ auquel appartiennent les propriétés abstraites suivantes :
\begin{enumerate}
\item[I.] La somme, la différence et le produit de deux nombres de $\mathfrak G$ sont encore des nombres de $\mathfrak G$.
\item[II.] Tout nombre réel ou complexe est quotient de deux nombres de $\mathfrak G$.
\item[III.] Tout nombre rationnel entier, respectivement algébrique entier, est un élément de $\mathfrak G$.
\item[IV.] Aucun nombre rationnel, respectivement algébrique, non entier n'appartient à $\mathfrak G$.\footnote{E. Zermelo, \emph{Über ganze transzendente Zahlen}, Math. Ann. 75, p. 434 (1914).}
\end{enumerate}

% unit: noether.p09.par.002
La construction repose sur le fait que le théorème du bon ordre permet de reconnaître l'existence d'une \emph{base algébrique} de tous les nombres, c'est-à-dire d'un \emph{système} $H$ \emph{de nombres} $\eta$ entre lesquels il n'existe aucune relation algébrique, mais au moyen desquels tous les autres nombres peuvent s'exprimer algébriquement. En raison de leur indépendance algébrique, ces nombres de base $\eta$ peuvent être regardés comme des indéterminées; le domaine recherché devient donc un domaine entier de fonctions rationnelles et algébriques des indéterminées $\eta$,\footnote{C'est pourquoi les considérations du présent article sont en partie parallèles à celles de l'article antérieur de l'autrice : \emph{Körper und Systeme rationaler Funktionen}, Math. Ann. 76, p. 161 (1915).} dont les coefficients appartiennent au corps $K$ de tous les nombres algébriques. Le domaine spécial $\mathfrak G_\eta$ construit par Zermelo est alors caractérisé par les propriétés de base suivantes :

% unit: noether.p09.par.003
$\mathfrak G_\eta$ contient :
\begin{enumerate}
\item[1)] tous les nombres de base $\eta$;
\item[2)] tous les polynômes en les $\eta$ à coefficients entiers;\footnote{Dans tout ce qui suit, \og{} entier \fg{} signifie que les coefficients appartiennent au domaine entier $[K]$ de tous les entiers algébriques. Pour la définition de $\mathfrak G_\eta$, il suffirait aussi, dans 1)--5), de ne considérer que des coefficients entiers rationnels; pour des domaines plus généraux, les propriétés de base deviendraient cependant moins simples.}
\item[3)] toutes les fonctions algébriques entières des $\eta$ à coefficients entiers.
\end{enumerate}
$\mathfrak G_\eta$ exclut :
\begin{enumerate}
\item[4)] toutes les fonctions rationnellement fractionnaires des $\eta$ à coefficients entiers;
\item[5)] toutes les fonctions algébriquement fractionnaires des $\eta$ à coefficients entiers.\footnote{Zermelo, loc. cit., \S{} 3.}
\end{enumerate}

% unit: noether.p09.par.004
Il est maintenant facile de voir, en comparant les exemples des \S{}\S{} 2 et 3, qu'il existe des domaines $\mathfrak G$ essentiellement différents du domaine de Zermelo $\mathfrak G_\eta$ : c'est-à-dire des domaines $\mathfrak G$ pour lesquels, sous aucun bon ordre, toutes les propriétés de base 1)--5) ne sont satisfaites, et qui par conséquent ne peuvent être produits par la construction de Zermelo sous aucun bon ordre. Autrement dit, il existe des domaines $\mathfrak G$ qui ne peuvent pas être appliqués biunivoquement et isomorphiquement sur $\mathfrak G_\eta$. Ainsi se pose la question des domaines les plus généraux d'entiers transcendants; elle sera complètement résolue dans ce qui suit.

% unit: noether.p09.par.005
À cette fin il faut d'abord déterminer lesquelles des propriétés de base sont des conséquences des propriétés abstraites de définition et lesquelles proviennent de la construction spéciale. On montrera (\S{} 1) que, pour tout domaine $\mathfrak G$ prescrit, il existe un bon ordre tel que les propriétés de base 1) et 2) soient satisfaites; ainsi tout domaine $\mathfrak G$ peut être construit au moyen d'une base algébrique de tous les nombres. Aucune des autres propriétés de base n'a besoin d'être satisfaite (\S{}\S{} 2 et 3). En particulier il s'ensuit qu'une autre propriété abstraite du domaine de Zermelo n'appartient pas à tous les domaines $\mathfrak G$ : la clôture algébrique-entière. Elle peut se formuler ainsi :
\begin{enumerate}
\item[V.] Tout nombre qui est entier sur le sous-anneau engendré par un nombre fini de nombres de $\mathfrak G$ appartient à $\mathfrak G$.
\end{enumerate}

% unit: noether.p09.par.006
Les domaines spéciaux pour lesquels V vaut en plus de I--IV seront appelés domaines $\mathfrak H$ d'entiers transcendants algébriquement-entiers. Pour les domaines les plus généraux $\mathfrak H$, on obtient les résultats simples suivants (\S{} 4) :
\begin{enumerate}
\item[a)] L'intersection de tous les domaines $\mathfrak H$ constructibles au moyen d'une même base $H$ est donnée par toutes les fonctions algébriques entières de la base, c'est-à-dire par le domaine de Zermelo $\mathfrak G_\eta$, qui est lui-même un domaine $\mathfrak H$. Cela donne aussi une caractérisation abstraite du domaine de Zermelo.
\item[b)] Tout domaine $\mathfrak H$ résulte de l'adjonction algébriquement entière\footnote{Explication de l'expression au \S{} 4.} d'un système arbitraire $\mathfrak S(\eta)$ à $\mathfrak G_\eta$, où $\mathfrak S(\eta)$ ne doit satisfaire qu'à la seule condition qu'aucun nombre algébriquement fractionnaire n'entre dans le domaine par cette adjonction.
\end{enumerate}

% unit: noether.p09.par.007
Les résultats relatifs aux domaines les plus généraux $\mathfrak G$, lorsqu'on prend pour fondement une base algébrique, sont moins simples. Ici l'intersection de tous les domaines $\mathfrak G$ n'est elle-même pas un domaine $\mathfrak G$, et la construction des domaines les plus généraux devient par conséquent plus difficile à embrasser (\S{} 5).

% unit: noether.p09.par.008
Afin d'obtenir des résultats analogues à ceux des domaines $\mathfrak H$, la base algébrique est remplacée par une \emph{base rationnelle} $\Theta$, c'est-à-dire par un \emph{système} $\Theta$ \emph{de nombres} $\vartheta$ qui permet d'exprimer tous les nombres rationnellement, tandis que, sous au moins un bon ordre, aucun nombre de base ne peut être exprimé rationnellement par un nombre fini de précédents. Cette base rationnelle $\Theta$ doit encore être soumise à la condition accessoire que le domaine entier $[\Theta]$, formé de tous les polynômes entiers en les $\vartheta$, ne contienne aucun nombre algébriquement fractionnaire. La construction d'une telle base avec condition accessoire est donnée au \S{} 7. Alors, pour les domaines les plus généraux $\mathfrak G$ --- que l'on pourrait aussi appeler domaines d'entiers transcendants rationnellement entiers ---, on obtient à nouveau les résultats simples suivants (\S{} 6) :
\begin{enumerate}
\item[a)] L'intersection de tous les domaines $\mathfrak G$ constructibles au moyen de la même base rationnelle $\Theta$ est donnée par toutes les fonctions rationnelles entières de la base, c'est-à-dire par le domaine $[\Theta]$, qui est lui-même un domaine $\mathfrak G$.
\item[b)] Tout domaine $\mathfrak G$ résulte de l'adjonction rationnellement entière d'un système arbitraire $\mathfrak S(\vartheta)$ à $[\Theta]$, où $\mathfrak S(\vartheta)$ ne doit satisfaire qu'à la seule condition qu'aucun nombre algébriquement fractionnaire n'entre dans le domaine par cette adjonction.
\end{enumerate}

% unit: noether.p09.par.009
Pour avoir une analogie complète avec les domaines $\mathfrak H$, il faudrait retrancher les nombres algébriques des propriétés de définition III et IV des domaines $\mathfrak G$. Avec cette modification de III et IV, la construction des éléments entiers au moyen d'une base rationnelle peut être transférée à un corps arbitraire (\S{} 10), tandis que la construction de Zermelo présuppose que le corps soit algébriquement clos.

% unit: noether.p09.par.010
Lorsque $H$ parcourt toutes les bases algébriques, respectivement lorsque $\Theta$ parcourt toutes les bases rationnelles satisfaisant à la condition accessoire, les résultats a) et b) donnent la totalité de tous les domaines $\mathfrak H$ et $\mathfrak G$. Si l'on réunit en une classe tous les domaines isomorphes, on peut alors extraire de cette totalité un sous-ensemble qui représente la totalité de toutes les classes --- au moins dénombrablement infinies d'après le \S{} 9 (\S{} 8).

% unit: noether.p09.sec.001
\subsection*{\S{} 1. Démonstration des propriétés de base 1) et 2) pour tous les domaines $\mathfrak G$.}

% unit: noether.p09.par.011
Soit $\mathfrak G$ un domaine arbitrairement prescrit d'entiers transcendants, de sorte qu'il possède les propriétés abstraites I--IV. En suivant le procédé appliqué par E. Zermelo au corps de tous les nombres complexes, choisissons dans $\mathfrak G$ une \emph{base algébrique} $H$ \emph{de tous les nombres de $\mathfrak G$}. Comme, d'après II, tout nombre réel ou complexe peut être représenté comme quotient de deux nombres de $\mathfrak G$, $H$ est en même temps une base algébrique de \emph{tous} les nombres. Ainsi, pour tout domaine $\mathfrak G$ prescrit, la base algébrique de tous les nombres peut être choisie de façon à appartenir à $\mathfrak G$; la propriété de base 1) est donc satisfaite. De plus, par III, $\mathfrak G$ contient le domaine entier $[K]$ de tous les entiers algébriques, et par conséquent, par I, aussi tous les polynômes en les éléments de base $\eta$ à coefficients dans $[K]$, c'est-à-dire tous les polynômes entiers en les $\eta$, qui forment le domaine entier $[H]$. La propriété de base 2) est donc elle aussi toujours satisfaite; autrement dit, tout domaine $\mathfrak G$ peut être construit au moyen d'une base algébrique $H$ de tous les nombres de telle manière que $\mathfrak G$ contienne le domaine entier $[H]$ comme diviseur.

% unit: noether.p09.par.012
\textbf{Remarque.} Néanmoins, en partant d'une base $H$, on peut naturellement construire un domaine $\mathfrak G$ pour lequel les propriétés de base 1) et 2) ne valent pas relativement à $H$ elle-même, mais valent relativement à une autre base $H'$. Par exemple, que $\mathfrak G'$ provienne du domaine de Zermelo $\mathfrak G_\eta$ en remplaçant, dans toute la construction, un certain élément de base $\eta$ par un polynôme $f(\eta)$. Alors $\mathfrak G'$ ne contient ni $\eta$ ni $[H]$; mais si, dans la base $H$, on remplace l'élément de base $\eta$ par $\xi=f(\eta)$ --- ce par quoi $H$ passe de nouveau à une base algébrique $H'$ ---, les propriétés de base 1) et 2) sont satisfaites pour $\mathfrak G'$ relativement à $H'$.

% unit: noether.p09.sec.002
\subsection*{\S{} 2. Exclusion des propriétés de base 4) et 5).}

% unit: noether.p09.par.013
Nous montrons maintenant que les domaines $\mathfrak G$ n'ont à satisfaire aucune des autres propriétés de base. D'abord 4) et 5) sont exclues en remplaçant, dans la construction de Zermelo, le domaine entier $[H]$ par un domaine de \og{} fonctionnelles \fg{} rationnelles entières.

% unit: noether.p09.par.014
D'après Weber,\footnote{H. Weber, \emph{Lehrbuch der Algebra}. Petite édition, \S{}\S{} 96 et 97.} on appelle \emph{fonctionnelle rationnelle} toute fonction rationnelle à coefficients rationnels écrite sous la représentation univoquement déterminée suivante :
% unit: noether.p09.eq.001
\[
 A(\eta_1\cdots \eta_t)
   = a\cdot\frac{E_1(\eta_1\cdots \eta_t)}{E_2(\eta_1\cdots \eta_t)},
\]
où $E_1,E_2$ sont des polynômes primitifs premiers entre eux à coefficients entiers rationnels, et où $a$ est un nombre rationnel positif, la valeur absolue de $A$. Si, en particulier, $a$ est un \emph{entier rationnel}, alors $A$ s'appelle une \emph{fonctionnelle rationnelle entière}. Comme cas particuliers, les fonctionnelles rationnelles entières comprennent les entiers rationnels et les polynômes à coefficients entiers rationnels; tandis que les nombres rationnellement fractionnaires et les polynômes à coefficients rationnellement fractionnaires doivent être rangés parmi les fonctionnelles rationnelles fractionnaires.

% unit: noether.p09.par.015
Soit maintenant le domaine $\mathfrak G$ constitué de la totalité $\mathfrak F$ de toutes les fonctionnelles rationnelles entières d'un nombre fini quelconque d'éléments de base $\eta$, et de toutes les fonctions qui sont entières sur $\mathfrak F$, c'est-à-dire qui satisfont à une équation dont le coefficient dominant est l'unité et dont les autres coefficients sont des fonctionnelles rationnelles entières.

% unit: noether.p09.par.016
La démonstration des propriétés I--IV pour $\mathfrak G$ est exactement parallèle à la démonstration correspondante chez Zermelo et ne sera qu'indiquée brièvement. D'abord II et III sont claires, puisque $\mathfrak G$ contient le domaine de Zermelo $\mathfrak G_\eta$ comme diviseur. Par le théorème de Gauss sur les polynômes à coefficients entiers rationnels, les sommes, différences et produits de fonctionnelles rationnelles entières sont encore des fonctionnelles rationnelles entières; donc $\mathfrak F$ forme un domaine entier, et par les arguments connus\footnote{Comparer par exemple Weber, \S{} 97, 8.} $\mathfrak G$ devient lui aussi un domaine entier; ainsi I est satisfaite. En outre, du théorème de Gauss étendu\footnote{Weber, \S{} 20, 5.} résultent les deux propositions auxiliaires suivantes : \og{} Si $z$ est entier sur $\mathfrak F$ en vertu d'une \emph{certaine} équation, alors il l'est aussi en vertu de l'équation \emph{irréductible} que $z$ satisfait relativement au corps de toutes les fonctionnelles rationnelles \fg{},\footnote{Comparer Weber, \S{} 97, 4.} et \og{} l'équation irréductible sur le corps de tous les nombres rationnels et satisfaite par un nombre algébrique reste irréductible sur le corps de toutes les fonctionnelles rationnelles \fg{}. Ainsi $\mathfrak G$ ne peut contenir aucun nombre algébrique fractionnaire; IV, et donc toutes les propriétés abstraites I--IV, sont vérifiées.

% unit: noether.p09.par.017
D'autre part, sous aucun bon ordre on ne peut choisir dans $\mathfrak G$ une base telle que $\mathfrak G$ exclue toutes les fonctions rationnellement et algébriquement fractionnaires des éléments de base. Soit en effet $z(\eta)$ une fonction quelconque du domaine qui doit être choisie comme élément de base, donc définie par une équation
% unit: noether.p09.eq.002
\[
 z^k+A_1(\eta)z^{k-1}+\cdots+A_k(\eta)=0,
\]
où
% unit: noether.p09.eq.003
\[
 A_i(\eta)=a_i\cdot\frac{E_1(\eta)}{E_2(\eta)}
\]
sont des fonctionnelles rationnelles entières. Alors la fonction $\frac{a_k}{z}$ appartient à $\mathfrak G$, de même que $\frac{a_k}{\sqrt z}$, $\frac{a_k}{\sqrt[3]z}$, etc.\footnote{Tout à fait de même, toute fonction algébrique des $\eta$ peut être transformée en élément du domaine fonctionnel $\mathfrak G$ par multiplication par un entier rationnel. Ce domaine possède ainsi la propriété que tout nombre complexe peut être représenté comme quotient d'un élément de $\mathfrak G$ et d'un entier rationnel, en analogie avec les nombres algébriques.} Les propriétés de base 4) et 5) sont donc exclues pour la totalité des domaines $\mathfrak G$.

% unit: noether.p09.par.018
\textbf{Remarque.} Un exemple au \S{} 3 montrera que $\mathfrak G$ peut, sous tout bon ordre, contenir des fonctionnelles rationnellement fractionnaires sans contenir un nombre rationnellement ou algébriquement fractionnaire.

% unit: noether.p09.sec.003
\subsection*{\S{} 3. Exclusion de la propriété de base 3).}

% unit: noether.p09.par.019
Pour exclure la propriété de base 3), nous utilisons un \emph{domaine de module} généralisé dans lequel les arguments des polynômes du domaine ne sont plus les indéterminées, mais toutes les fonctions algébriques entières des indéterminées,\footnote{Par \og{} fonctions algébriques entières \fg{} on entend toujours des fonctions à coefficients \emph{entiers}; c'est-à-dire des fonctions satisfaisant à une équation dont le coefficient dominant est l'unité et dont les autres coefficients sont des polynômes en les $\eta$ à coefficients entiers algébriques, donc des polynômes de $[H]$. En particulier, tous les polynômes de $[H]$ appartiennent aux fonctions algébriques entières.} tandis que les indéterminées elles-mêmes jouent le rôle d'une base du module.

% unit: noether.p09.par.020
Que $\mathfrak G$ consiste en \emph{tous les éléments}
% unit: noether.p09.eq.004
\begin{equation}
 z=\alpha+g_1(\eta)\cdot\eta_1+g_2(\eta)\cdot\eta_2+\cdots+g_t(\eta)\cdot\eta_t,
\tag{1}
\end{equation}
\emph{où $\alpha$ parcourt tous les entiers algébriques, $\eta_1\cdots\eta_t$ tous les éléments de base, et où, pour chaque combinaison fixe $\eta_1\cdots\eta_t$, les fonctions $g_1(\eta)\cdots g_t(\eta)$ parcourent individuellement toutes les fonctions algébriques entières des $\eta$.}\footnote{La propriété de module appartient aux éléments $z-\alpha$.}

% unit: noether.p09.par.021
Il est facile de vérifier que les propriétés I--IV valent pour ce domaine de module. D'abord I vaut, puisque la somme, la différence et le produit de deux éléments $z$ prennent de nouveau la forme (1). Le domaine $\mathfrak G$ contient aussi tous les éléments de base $\eta$, et en outre tous les éléments $\eta\cdot g(\eta)$, où $g(\eta)$ parcourt toutes les fonctions algébriques entières des $\eta$. Par conséquent toutes les fonctions algébriques entières des $\eta$, et donc toutes les fonctions algébriques des $\eta$ en général, peuvent être représentées comme quotients de deux éléments de $\mathfrak G$, ce qui prouve II. Dans la représentation (1), tous les entiers algébriques sont en particulier inclus --- pour $g_1=0\cdots g_t=0$ ---; mais $z$ ne peut pas être égal à un nombre algébrique fractionnaire. Car si $z$ représente un nombre algébrique $\beta$, alors de
% unit: noether.p09.eq.005
\[
 \beta=\alpha+g_1(\eta)\eta_1+\cdots+g_t(\eta)\eta_t
\]
identiquement en les $\eta$, il résulte nécessairement $\beta=\alpha$. Cela se voit par la spécialisation
% unit: noether.p09.eq.006
\[
 \eta_1=0\cdots \eta_t=0,
\]
sous laquelle les multiplicateurs $g_1(\eta)\cdots g_t(\eta)$ restent finis comme fonctions algébriques entières, puisqu'ils passent en des fonctions ou nombres algébriques entiers.\footnote{Cette preuve plus détaillée de IV est donnée à cause de l'exemple élargi à la fin du paragraphe. Ici suffirait la remarque que, par construction, $z$ est une fonction algébrique entière.} Ainsi les conditions I--IV sont satisfaites pour $\mathfrak G$.

% unit: noether.p09.par.022
D'autre part, sous aucun bon ordre on ne peut choisir dans $\mathfrak G$ une base telle que toutes les fonctions algébriques entières des éléments de base appartiennent à $\mathfrak G$. Soit $z(\eta)$ une fonction quelconque du domaine qui doit être choisie comme élément de base --- et qui doit donc effectivement contenir les indéterminées $\eta$. Nous montrerons qu'à partir d'une certaine valeur finie de $\nu$, dépendant de $z$, les fonctions
% unit: noether.p09.eq.007
\[
 \xi=\sqrt[\nu]{z-\alpha}
\]
ne peuvent plus appartenir à $\mathfrak G$.

% unit: noether.p09.par.023
Supposons en effet que $\xi$ appartienne à $\mathfrak G$. Alors, par (1), il est de la forme
% unit: noether.p09.eq.008
\[
 \xi=\gamma+h_1(\eta)\eta_{i_1}+h_2(\eta)\eta_{i_2}
      +\cdots+h_\tau(\eta)\eta_{i_\tau},
\]
où $h_1(\eta)\cdots h_\tau(\eta)$ sont encore des fonctions algébriques entières des $\eta$, et où $\eta_{i_1}\cdots\eta_{i_\tau}$ sont des éléments de base quelconques, qui peuvent en particulier coïncider aussi avec $\eta_1\cdots\eta_t$. Il faut d'abord montrer que l'entier algébrique $\gamma$ est nul. Comme $h_1(\eta)\cdots h_\tau(\eta)$ sont des fonctions algébriques entières, $\xi$ devient égal à $\gamma$ pour $\eta_{i_1}=0\cdots\eta_{i_\tau}=0$, avec des valeurs arbitraires des autres $\eta$; en particulier, $\xi$ est alors égal à $\gamma$ pour
% unit: noether.p09.eq.009
\[
 \eta_{i_1}=0\cdots \eta_{i_\tau}=0; \qquad
 \eta_1=0\cdots \eta_t=0.
\]
Mais sous cette spécialisation la fonction $(z-\alpha)$ s'annule par (1), et donc $\xi$ aussi; par conséquent $\gamma=0$, et l'on a
% unit: noether.p09.eq.010
\[
 \xi=h_1(\eta)\eta_{i_1}+h_2(\eta)\eta_{i_2}
      +\cdots+h_\tau(\eta)\eta_{i_\tau}.
\]
En élevant à la puissance $\nu$-ième, on obtient
% unit: noether.p09.eq.011
\[
 \xi^\nu=z-\alpha=F_\nu(\eta),
\]
où $F_\nu(\eta)$ désigne une forme homogène, non identiquement nulle, de dimension $\nu$-ième en $\eta_{i_1}\cdots\eta_{i_\tau}$, dont les coefficients sont des fonctions algébriques entières des $\eta$. Or $z-\alpha$, qui par (1) est une fonction algébrique entière des $\eta$, satisfait à une équation irréductible relativement à $[H]$ :
% unit: noether.p09.eq.012
\[
 (z-\alpha)^\chi+B(\eta)(z-\alpha)^{\chi-1}+\cdots+C(\eta)=0,
\]
où le dernier coefficient $C(\eta)$ --- produit de $(z-\alpha)$ avec ses conjugués --- est un polynôme; soit $\lambda$ son degré. D'autre part, de $z-\alpha=F_\nu(\eta)$, en multipliant $F_\nu(\eta)$ par les conjugués correspondants, on obtient
% unit: noether.p09.eq.013
\[
 C(\eta)=G_{\nu\chi}(\eta),
\]
où $G_{\nu\chi}(\eta)$ désigne une forme homogène de dimension $\nu\chi$-ième en $\eta_{i_1}\cdots\eta_{i_\tau}$, dont les coefficients sont des polynômes en les $\eta$. Ainsi $G_{\nu\chi}$ est d'un degré au moins $\nu\chi$ en les $\eta$, ou $\lambda\geq\nu\chi$; c'est-à-dire que les fonctions $\sqrt[\nu]{z-\alpha}$ pour lesquelles $\nu>\lambda/\chi$ n'appartiennent pas à $\mathfrak G$. La propriété de base 3) est donc exclue pour la totalité des domaines $\mathfrak G$.

% unit: noether.p09.par.024
\textbf{Remarque.} Une légère généralisation du domaine qui vient d'être construit conduit à un domaine $\mathfrak G$ qui, sous tout bon ordre, contient aussi des fonctionnelles fractionnaires. Que $\mathfrak G$ consiste en tous les éléments
% unit: noether.p09.eq.014
\[
 z=\alpha+\gamma_1(\eta)\eta_1+\gamma_2(\eta)\eta_2
      +\cdots+\gamma_t(\eta)\eta_t,
\]
où maintenant $\gamma(\eta)$ parcourt toutes les fonctions algébriquement entières, mais non nécessairement entières, des $\eta$. La preuve des propriétés I--IV reste la même que ci-dessus. Mais puisque $\mathfrak G$, outre toute fonction $z$, contient aussi $a\cdot(z-\alpha)$, où $a$ désigne un nombre rationnel ou algébriquement fractionnaire quelconque, des fonctionnelles fractionnaires apparaissent également sous tout bon ordre.

% unit: noether.p09.sec.004
\subsection*{\S{} 4. Les domaines les plus généraux d'entiers transcendants algébriquement-entiers.}

% unit: noether.p09.par.025
Pour formuler plus commodément ce qui suit, nous dirons qu'un élément $z$ est \emph{entier sur $\mathfrak T$} lorsque $z$ satisfait à une équation dont le coefficient dominant est l'unité et dont les autres coefficients sont des polynômes entiers en les éléments de $\mathfrak T$, où $\mathfrak T$ désigne un domaine quelconque prescrit.\footnote{Pour un domaine entier $\mathfrak T$, le coefficient dominant est l'unité et les autres coefficients sont des éléments de $\mathfrak T$.}

% unit: noether.p09.par.026
Un domaine $\mathfrak T$ s'appelle \emph{intégralement clos} s'il satisfait à la condition, comparer l'introduction :
\begin{enumerate}
\item[V.] Tout élément entier sur $\mathfrak T$ appartient à $\mathfrak T$.
\end{enumerate}

% unit: noether.p09.par.027
Nous montrons d'abord que la propriété abstraite V appartient au domaine de Zermelo $\mathfrak G_\eta$. Soit un élément $z$ entier sur $\mathfrak G_\eta$ en vertu d'une équation $f(z)=0$. Alors la multiplication de $f(z)$ par ses conjuguées relativement à $[H]$ montre que $z$ est aussi entier sur $[H]$, et appartient donc à $\mathfrak G_\eta$. La clôture intégrale du domaine fonctionnel considéré au \S{} 2 se démontre de la même manière.

% unit: noether.p09.par.028
Que la propriété V n'appartienne pas à tout domaine $\mathfrak G$ est montré par le domaine de module considéré au \S{} 3, qui ne possède précisément pas la propriété de base 3). Plus exactement, le \S{} 3 a montré que le domaine de module n'est intégralement clos relativement à aucun élément transcendant du domaine, tandis que, par III et IV, c'est bien entendu le cas relativement à tout élément algébrique du domaine.

% unit: noether.p09.par.029
Cela nous conduit d'abord à choisir, dans la totalité des domaines $\mathfrak G$, ceux qui possèdent aussi la propriété abstraite V; ceux-ci seront appelés \og{} domaines $\mathfrak H$ d'entiers transcendants algébriquement-entiers \fg{}.

% unit: noether.p09.par.030
Soit $\mathfrak H$ un tel domaine, et soit $H$ une base algébrique de tous les nombres de $\mathfrak H$; d'après II, $H$ est en même temps une base algébrique de tous les nombres. D'après III, I et V, $\mathfrak H$ contient, outre $H$, toutes les fonctions algébriques entières des $\eta$, c'est-à-dire le domaine de Zermelo $\mathfrak G_\eta$. Ainsi, en analogie avec le \S{} 1 : tout domaine $\mathfrak H$ d'entiers transcendants algébriquement-entiers peut être construit au moyen d'une base algébrique de tous les nombres de telle manière que $\mathfrak H$ contienne le domaine de Zermelo $\mathfrak G_\eta$ comme diviseur.

% unit: noether.p09.par.031
Soit maintenant $H$ une base algébrique quelconque de tous les nombres, et que $\mathfrak M(\mathfrak H)$ désigne la totalité de tous les domaines $\mathfrak H$ contenant $H$. Tout domaine $\mathfrak H$ de $\mathfrak M(\mathfrak H)$ contient, comme on vient de le prouver, le domaine de Zermelo $\mathfrak G_\eta$ comme diviseur; et puisque $\mathfrak G_\eta$ lui-même est un domaine $\mathfrak H$, $\mathfrak G_\eta$ est le plus grand commun diviseur, ou l'intersection, de tous les domaines $\mathfrak H$ de $\mathfrak M(\mathfrak H)$. Ainsi $\mathfrak G_\eta$ est défini abstraitement. Il en résulte aussi directement que $\mathfrak M(\mathfrak H)$ est fermé pour l'intersection : c'est-à-dire que l'intersection de tous les domaines $\mathfrak H$ d'un sous-système quelconque de $\mathfrak M(\mathfrak H)$ appartient encore à $\mathfrak M(\mathfrak H)$. En effet, comme intersection, I et V sont satisfaites; II et III résultent de ce que $\mathfrak G_\eta$ est diviseur; IV résulte de ce que l'intersection est diviseur de $\mathfrak H$.

% unit: noether.p09.par.032
Comme le montre l'exemple du domaine fonctionnel au \S{} 2, les domaines $\mathfrak H$ contiennent en général d'autres fonctions encore que celles de $\mathfrak G_\eta$. Soit $\psi(\eta)$ une telle fonction rationnelle ou algébrique des $\eta$. Alors, par I, $\mathfrak H$ contient aussi toutes les combinaisons rationnelles entières de $\psi$ et, chaque fois, d'un nombre fini de fonctions de $\mathfrak G_\eta$. Le domaine entier ainsi obtenu --- constitué de tous les polynômes en $\psi$ à coefficients dans $\mathfrak G_\eta$ --- sera noté $[\mathfrak G_\eta,\psi]$, et le procédé qui y conduit sera appelé \og{} adjonction polynomiale de $\psi$ à $\mathfrak G_\eta$ \fg{}. Par V, $\mathfrak H$ contient en outre toutes les fonctions entières sur $[\mathfrak G_\eta,\psi]$; le domaine ainsi obtenu sera noté $\{\mathfrak G_\eta,\psi\}$, et le procédé qui y conduit \og{} adjonction intégrale de $\psi$ à $\mathfrak G_\eta$ \fg{}. Le domaine $\{\mathfrak G_\eta,\psi\}$ se compose de toutes les fonctions algébriques entières de $\psi$ à coefficients dans $\mathfrak G_\eta$ et est donc, par les arguments connus, un domaine entier.\footnote{Comparer par exemple Weber, \S{} 97, 6 et 7.}

% unit: noether.p09.par.033
Nous montrons que $\{\mathfrak G_\eta,\psi\}$ est aussi intégralement clos, c'est-à-dire satisfait à la condition V. En effet, que $z$ soit entier sur $\{\mathfrak G_\eta,\psi\}$ en vertu d'une équation
% unit: noether.p09.eq.015
\[
 f(z;a\cdots c)=z^\chi+a z^{\chi-1}+\cdots+c=0,
\]
où $a\cdots c$, comme éléments de $\{\mathfrak G_\eta,\psi\}$, satisfont à des équations à coefficients dans $[\mathfrak G_\eta,\psi]$ :
% unit: noether.p09.eq.016
\[
\begin{gathered}
 a^\lambda+A_1a^{\lambda-1}+\cdots+A_\lambda=0,\\
 \vdots\\
 c^\mu+C_1c^{\mu-1}+\cdots+C_\mu=0.
\end{gathered}
\]
Désignons leurs racines par $a,a'\cdots a^{(\lambda-1)}\cdots c,c'\cdots c^{(\mu-1)}$. Si l'on forme maintenant le produit sur toutes les combinaisons des racines,
% unit: noether.p09.eq.017
\[
 g(z)=f(z;a\cdots c)f(z;a'\cdots c')\cdots
      f(z;a^{(\lambda-1)}\cdots c^{(\mu-1)}),
\]
alors $g(z)$ a pour coefficient dominant l'unité et pour autres coefficients des éléments de $[\mathfrak G_\eta,\psi]$. Par conséquent $z$ appartient à $\{\mathfrak G_\eta,\psi\}$.\footnote{Ainsi, en définissant l'intégralité algébrique au moyen d'une équation \emph{quelconque}, la notion d'irréductibilité relativement à un corps de base devient inutile pour la preuve de I et V. Comparer aussi le \S{} 2, où, seulement pour la preuve de IV, on a besoin de la proposition auxiliaire qui fait passer de la définition de l'intégralité algébrique par une équation arbitraire à l'équation irréductible. Comme cette proposition auxiliaire ne vaut plus en général dans le cas présent, IV doit être imposée à la fonction $\psi$ comme condition spéciale.}

% unit: noether.p09.par.034
Ainsi le domaine $\{\mathfrak G_\eta,\psi\}$ possède les propriétés I et V; et, puisque $\mathfrak G_\eta$ est diviseur du domaine, il possède aussi II et III. Que IV soit satisfaite dépend du choix de $\psi$; par exemple IV ne l'est pas pour $\psi=\frac{1}{n\cdot\eta}$, car alors $\frac1n$ appartient au domaine. Elle est toutefois satisfaite, d'après le \S{} 2, pour $\psi=\frac1\eta$, ou encore pour $\psi=\frac{\eta}{n}$. Dans ce dernier cas, en effet, si une fonction de $\{\mathfrak G_\eta,\psi\}$ représente un nombre algébrique, sa valeur reste inchangée pour $\eta=0$; elle passe ainsi en une fonction de $\mathfrak G_\eta$, et par conséquent en un entier algébrique.

% unit: noether.p09.par.035
Ainsi $\{\mathfrak G_\eta,\psi\}$ devient un domaine $\mathfrak H$ dès que l'on impose à $\psi$ la condition qu'aucun nombre algébriquement fractionnaire n'entre dans le domaine par l'adjonction algébriquement entière.

% unit: noether.p09.par.036
De manière tout à fait analogue, on définit l'adjonction polynomiale, respectivement intégrale, d'un système arbitraire $\mathfrak S$ de fonctions rationnelles ou algébriques des $\eta$ à $\mathfrak G_\eta$. L'adjonction polynomiale conduit au domaine entier $[\mathfrak G_\eta,\mathfrak S]$, constitué de tous les polynômes entiers en un nombre fini d'éléments de $\mathfrak G_\eta$ et de $\mathfrak S$ à la fois. L'adjonction intégrale donne le domaine $\{\mathfrak G_\eta,\mathfrak S\}$ constitué de tous les éléments entiers sur $[\mathfrak G_\eta,\mathfrak S]$. Exactement comme ci-dessus, on montre que le domaine entier $\{\mathfrak G_\eta,\mathfrak S\}$ est intégralement clos et satisfait aussi II et III. Ainsi, comme ci-dessus :

% unit: noether.p09.par.037
$\{\mathfrak G_\eta,\mathfrak S\}$ devient un domaine $\mathfrak H$ dès que l'on impose à $\mathfrak S$ la condition qu'aucun nombre algébriquement fractionnaire n'entre dans le domaine par l'adjonction intégrale --- c'est-à-dire dès que $\mathfrak S$ devient un système admissible.\footnote{Plus généralement, on montre de même : si $\mathfrak T$ est un domaine quelconque et si $\{\mathfrak T\}$ désigne la totalité des éléments entiers sur $\mathfrak T$, alors $\{\mathfrak T\}$ est intégralement clos.}

% unit: noether.p09.par.038
Il est maintenant important que la réciproque soit également vraie, de sorte que les domaines $\{\mathfrak G_\eta,\mathfrak S\}$ épuisent effectivement tous les domaines $\mathfrak H$ de $\mathfrak M(\mathfrak H)$ lorsque $\mathfrak S$ parcourt tous les systèmes admissibles. En effet, soit $\mathfrak H$ un domaine quelconque de $\mathfrak M(\mathfrak H)$, et désignons par $\mathfrak L=\mathfrak H-\mathfrak G_\eta$ le système constitué de toutes les fonctions de $\mathfrak H$ n'appartenant pas à $\mathfrak G_\eta$. Par V, $\mathfrak H$ contient le domaine $\{\mathfrak G_\eta,\mathfrak L\}$ comme diviseur; mais $\{\mathfrak G_\eta,\mathfrak L\}$ est aussi divisible par $\mathfrak H$, car il contient $\mathfrak G_\eta$ et $\mathfrak L$, et, puisque ceux-ci n'ont aucun élément commun, il contient aussi $\mathfrak H$. Donc $\mathfrak H=\{\mathfrak G_\eta,\mathfrak L\}$; et puisque IV est satisfaite pour $\mathfrak H$, $\mathfrak L$ est naturellement un système admissible.\footnote{L'adjonction d'un sous-système de $\mathfrak L$ peut déjà suffire; par exemple le domaine fonctionnel du \S{} 2 provient de l'adjonction algébriquement entière de toutes les fonctionnelles rationnelles entières --- à l'exception des polynômes --- à $\mathfrak G_\eta$.}

% unit: noether.p09.par.039
Lorsque $H$ parcourt toutes les bases algébriques possibles, $\mathfrak M(\mathfrak H)$ parcourt, comme on l'a prouvé au début, la totalité de tous les domaines $\mathfrak H$. Par conséquent la question des domaines les plus généraux $\mathfrak H$ d'entiers transcendants algébriquement-entiers est complètement résolue par les résultats de ce paragraphe :
\begin{enumerate}
\item[a)] L'intersection de tous les domaines $\mathfrak H$ de $\mathfrak M(\mathfrak H)$ est donnée par le domaine de Zermelo $\mathfrak G_\eta$, qui est lui-même un domaine $\mathfrak H$; $\mathfrak M(\mathfrak H)$ est fermé relativement à la formation des intersections.
\item[b)] Tout domaine $\mathfrak H$ de $\mathfrak M(\mathfrak H)$ provient de l'adjonction intégrale d'un système admissible $\mathfrak S$ à $\mathfrak G_\eta$. Lorsque $H$ parcourt toutes les bases algébriques possibles, $\mathfrak M(\mathfrak H)$ parcourt la totalité de tous les domaines $\mathfrak H$.\footnote{Les systèmes admissibles $\mathfrak S$ peuvent être trouvés comme suit. On place toutes les fonctions algébriquement fractionnaires des $\eta$ sous un bon ordre $\Omega$, et l'on prend dans $\mathfrak S$ toutes les fonctions sauf celles dont l'adjonction intégrale à $\mathfrak G_\eta$ et aux fonctions admises auparavant engendre un nombre algébriquement fractionnaire. Ce procédé peut être interrompu en un lieu arbitraire. Lorsque $\Omega$ parcourt tous les bons ordres, et que le procédé est interrompu en tout lieu arbitraire pour chaque $\Omega$ fixé, tous les systèmes admissibles $\mathfrak S$ sont ainsi épuisés.}
\end{enumerate}
